%==============================================================================
% Anexos
%==============================================================================

\section{Diagrama de arquitectura del sistema}

El Anexo~A reune los diagramas de la arquitectura propuesta en notacion UML. Incluye:

\begin{itemize}
  \item \textbf{Diagrama de despliegue}: los cinco componentes principales (N8N, FastAPI, PostgreSQL, Gemini~2.5 Flash, Twilio) y sus relaciones de comunicacion cifrada sobre TLS~1.3.
  \item \textbf{Diagrama de secuencia}: flujo extremo a extremo de un incidente desde su recepcion en el canal de correo electronico hasta la confirmacion al usuario, incluyendo las interacciones con el clasificador hibrido y la base de datos.
  \item \textbf{Diagrama de componentes}: organizacion interna del modulo Python por capas (routes, services, repositories, models, classifiers), con sus dependencias y relaciones.
\end{itemize}

Los diagramas se mantienen en formato editable (\texttt{.drawio}) en el repositorio del proyecto para facilitar su actualizacion futura. La version en alta resolucion se encuentra disponible en el directorio \texttt{docs/} del repositorio \parencite{Gomez2026_repositorio}.

\section{Repositorio de codigo fuente}

El codigo fuente del modulo Python, los scripts de migracion de la base de datos PostgreSQL gestionados mediante Alembic, los archivos de configuracion del despliegue en contenedores Docker y los flujos de integracion continua sobre GitHub Actions se encuentran disponibles en el repositorio publico del proyecto \parencite{Gomez2026_repositorio}. El commit de referencia para la version entregada al jurado corresponde al tag \texttt{v1.0.0}.

La estructura del repositorio sigue el patron estandar con directorios separados para el modulo de procesamiento (\texttt{Gestion\_Incidentes/}), la definicion del flujo de orquestacion (\texttt{Automatizacion\_Mesa\_de\_Ayuda.json}), las migraciones de base de datos (\texttt{alembic/}), las pruebas automatizadas (\texttt{tests/}), los artefactos de evaluacion (\texttt{evaluation/}) y la documentacion operativa (\texttt{docs/}). El archivo \texttt{README.md} incluye instrucciones detalladas para el despliegue local en menos de 15 minutos.

\section{Esquema de base de datos}

El esquema completo de la base de datos PostgreSQL se documenta mediante scripts SQL versionados que definen las cinco tablas descritas en la Tabla~4 del Capitulo~5, las restricciones de integridad referencial (claves foraneas con eliminacion en cascada condicional), los indices secundarios sobre los atributos de busqueda frecuente (fecha de creacion, sector responsable y estado actual del ticket), y las funciones almacenadas que generan los identificadores secuenciales con prefijo configurable. Las migraciones incrementales se gestionan mediante Alembic, lo que permite reproducir el estado de la base en cualquier punto historico del desarrollo y revertir cambios cuando resulte necesario.

La migracion inicial incluye la creacion de las tablas de catalogo (\texttt{sector}, \texttt{estado}, \texttt{canal\_origen}) con sus valores semilla, garantizando que un despliegue nuevo del sistema comience con los datos de referencia correctos. Las migraciones subsiguientes documentan la evolucion del esquema durante las seis iteraciones de desarrollo descritas en el Capitulo~4.

\section{Especificacion OpenAPI de la interfaz REST}

La especificacion completa de la interfaz REST se genera automaticamente mediante FastAPI conforme a la version~3.1 de la especificacion OpenAPI \parencite{Tiangolo2024_fastapi}. El documento resultante describe los puntos de entrada listados en la Tabla~5 del Capitulo~5, los esquemas de los cuerpos de solicitud y respuesta validados mediante Pydantic~v2, los codigos de estado posibles para cada metodo HTTP, y ejemplos representativos para cada caso de uso.

La interfaz se encuentra disponible adicionalmente en formato interactivo a traves del componente Swagger UI integrado en el servidor, accesible bajo la ruta \texttt{/docs} durante la ejecucion, y en formato estatico como archivo \texttt{openapi.json} en el directorio \texttt{docs/} del repositorio. La sincronizacion entre el codigo y la especificacion estatica se verifica automaticamente en el pipeline de integracion continua.

\section{Configuracion del flujo N8N}

La exportacion del flujo de orquestacion en formato JSON incluye la totalidad de los 12 nodos configurados, sus parametros operativos, las credenciales referenciadas mediante identificadores opacos sin exposicion de valores, y los conectores entre nodos. El archivo \texttt{Automatizacion\_Mesa\_de\_Ayuda.json} en la raiz del repositorio contiene esta exportacion. La importacion de este artefacto en una instancia limpia de N8N~1.62 reproduce integramente el comportamiento operativo del sistema, sujeta unicamente a la provision de las credenciales correspondientes a los servicios externos integrados (servidor IMAP de correo, cuenta Twilio y clave de API de Gemini) mediante el sistema de credenciales nativo de la plataforma.

\section{Corpus de validacion}

El corpus de 200 incidentes utilizado para la validacion experimental se conserva en formato CSV en el directorio \texttt{evaluation/data/} del repositorio, con columnas para el identificador anonimo del caso, la descripcion cruda pseudonimizada, el canal de origen, la categoria asignada por consenso humano, la categoria predicha por el clasificador, el tiempo de registro medido en cada flujo y la confianza asociada a la prediccion. Se incluyen ademas las planillas de cronometraje del flujo manual y los registros automaticos generados por el flujo automatizado.

Por razones de privacidad y en cumplimiento de la Ley~25.326 \parencite{Ley25326_2000}, las descripciones textuales han sido pseudonimizadas previamente a su inclusion en el repositorio publico, reemplazando nombres propios, direcciones de correo y otros identificadores personales por etiquetas genericas.

\section{Documentacion operativa}

La documentacion operativa del sistema describe los procedimientos de despliegue inicial, configuracion recurrente, respaldo y restauracion, monitoreo activo y respuesta ante incidentes operativos. Se incluyen instrucciones detalladas para los entornos de desarrollo, preproduccion y produccion, asi como una guia de resolucion de incidencias frecuentes orientada al personal de operaciones de la organizacion adoptante. La documentacion se mantiene actualizada en el repositorio del proyecto y se versiona conjuntamente con el codigo fuente, garantizando coherencia entre la version desplegada y los procedimientos vigentes en cada momento del ciclo de vida del sistema.

\newpage
